% Bagian: conditional-logics
% Bab: minimal-change-semantics
% Subbagian: true-false

\documentclass[../../../include/open-logic-section]{subfiles}

\begin{document}

\olfileid[id]{con}{min}{tf}

\olsection{Kebenaran dan Kepalsuan Kontrafaktual}

Suatu kontrafaktual $!A \cif !B$ benar (secara takvakum) jika semua
dunia-$!A$ yang paling dekat merupakan dunia-$!B$, sebagaimana digambarkan
dalam \olref{fig:true}.
\begin{figure}
\begin{center}
\begin{tikzpicture}[scale=.7]\small
  \spheresystem{5}
  \draw (0,0) node {$w$};
  \propositionintersect{0}{4}{40}{3.5}
  {\tikzset{proposition/.append={smooth,tension=1.4}}
  \proposition[shift={(1.6,-1)}]{90}{1}{30}{4}}
  \path (1.5,3) node[below] {$\formula{B}$};
  \path (3,0) node {$\formula{A}$};
\end{tikzpicture}
\caption{Kontrafaktual yang benar secara takvakum}
\ollabel{fig:true}
\end{center}
\end{figure}
Suatu kontrafaktual juga benar pada $w$ jika sistem bola di sekitar~$w$
sama sekali tidak mempunyai bola yang memuat dunia-$!A$. Dalam hal itu,
kontrafaktual tersebut benar secara \emph{vakum} (lihat
\olref{fig:vacuous}).
\begin{figure}
\begin{center}
\begin{tikzpicture}[scale=.7]\small
  \spheresystem{5}
  \draw (0,0) node {$w$};
  \proposition{0}{6.5}{40}{3.5}
  {\tikzset{proposition/.append={smooth,tension=1.4}}
  \proposition[shift={(1.2,-1)}]{90}{1}{30}{4}}
  \path (1.5,3) node[below] {$\formula{B}$};
  \spherepos{0}{7}{node {$\formula{A}$}}
\end{tikzpicture}
\caption{Kontrafaktual yang benar secara vakum}
\ollabel{fig:vacuous}
\end{center}
\end{figure}

Kontrafaktual dapat salah dengan dua cara. Cara pertama terjadi jika tidak
semua dunia-$!A$ yang paling dekat merupakan dunia-$!B$, tetapi beberapa di
antaranya memang demikian. Dalam hal ini, $!A \cif \lnot !B$ juga salah
(lihat \olref{fig:false}).
\begin{figure}
\begin{center}
\begin{tikzpicture}[scale=.7]\small
  \spheresystem{5}
  \draw (0,0) node {$w$};
  \propositionintersect{0}{4}{40}{3.5}
  {\tikzset{proposition/.append={smooth,tension=1.4}}
  \proposition[shift={(1.6,0)}]{90}{1}{30}{3}}
  \path (1.5,3) node[below] {$\formula{B}$};
  \path (3,0) node {$\formula{A}$};
\end{tikzpicture}
\caption{Kontrafaktual salah, lawannya salah}
\ollabel{fig:false}
\end{center}
\end{figure}
Jika dunia-dunia-$!A$ yang paling dekat sama sekali tidak beririsan dengan
dunia-dunia-$!B$, maka $!A \cif !B$ salah. Namun, dalam hal ini semua
dunia-$!A$ yang paling dekat merupakan dunia-$\lnot !B$, sehingga $!A \cif
\lnot !B$ benar (lihat \olref{fig:false-opposite}).
\begin{figure}
\begin{center}
\begin{tikzpicture}[scale=.7]\small
  \spheresystem{5}
  \draw (0,0) node {$w$};
  \propositionintersect{0}{4}{40}{3.5}
  {\tikzset{proposition/.append={smooth,tension=1.4}}
  \proposition[shift={(-1,0)}]{90}{1}{30}{3}}
  \path (-1,3) node[below] {$\formula{B}$};
  \path (3,0) node {$\formula{A}$};
\end{tikzpicture}
\caption{Kontrafaktual salah, lawannya benar}
\ollabel{fig:false-opposite}
\end{center}
\end{figure}

Berbeda dari kondisional ketat, kontrafaktual dapat bersifat kontingen.
Perhatikan model bola dalam \olref{fig:contingent}. Semua dunia-$!A$ yang
paling dekat dengan~$u$ merupakan dunia-$!B$, sehingga $\mSat{M}{!A \cif
  !B}[u]$. Namun, terdapat dunia-dunia-$!A$ yang paling dekat dengan~$v$ yang
bukan dunia-$!B$, sehingga $\mSat/{M}{!A \cif !B}[v]$.
\begin{figure}
\begin{center}
\begin{tikzpicture}[scale=.6]\small
  \spheresystem{7}
  \spheresystem[shift={(0:3.1)},dashed]{4}
  \propositionintersect{45}{4}{40}{4.5}
  \begin{scope}
    \clip \propositionplot{45}{4}{40}{4.5} ;
    \spherelayer[shift={(0:3.1)}]{3}
  \end{scope}
  \proposition[shift={(1.4,-.5)}]{90}{1}{35}{5}
  \path (0,0) node {$u$};
  \path (0:3.1) node {$v$};
  \spherepos{35}{9}{node {$\formula{A}$}}
  \spherepos{80}{9}{node {$\formula{B}$}}
\end{tikzpicture}
\end{center}
\caption{Kontrafaktual kontingen}
\ollabel{fig:contingent}
\end{figure}
\end{document}
